\documentclass[a4paper, 11pt]{article} % das Papierformat zuerst
\usepackage[utf8]{inputenc}
\usepackage{geometry}
\geometry{a4paper, top=25mm, left=35mm, right=25mm, bottom=25mm}
\usepackage{graphicx}
\usepackage{color}
\usepackage{epstopdf}
\usepackage[T1]{fontenc}
\usepackage{setspace}
\usepackage{tabularx}
\usepackage{blindtext}
\usepackage[ngerman]{babel} %deutsche Silbentrennung
\usepackage{titlesec}
\usepackage{enumitem} 
\usepackage{hyperref}
\usepackage{textcomp}



\definecolor{fsi}{RGB}{0,81,150}

% VARIABLEN
\newcommand{\protokoller}{Dominik Petersdorf}
\newcommand{\dateOfMeeting}{21. Mai 2015}
\newcommand{\TeXer}{Dominik Petersdorf}
\newcommand{\fsiPresident}{Georg Schäfer}


\begin{document}
%deckblatt start

\doublespacing
\thispagestyle{empty}

\begin{center}
\includegraphics[width=10.0cm]{../logo_faculty_computer_science.eps}

\vspace*{\fill}
\begin{singlespace}
	{\LARGE \textbf{Protokoll der Vollversammlung der Fachschaft Informatik 
			\singlespacing
			vom \dateOfMeeting}}\\
			Fachschaftsrat Informatik\\
			Trier University of Applied Sciences\\
			\vspace{2.5cm}
			\textit{
			Protokollführer: \textbf{\protokoller} \\
			\LaTeX - Umsetzung von \TeXer\\
			am \today
	}
\end{singlespace}

\vfill
\end{center}

\hspace*{-35cm}
\textcolor{fsi}{\rule{64.9cm}{15pt}}
\pagebreak
 
\setcounter{tocdepth}{2}
\tableofcontents 
\pagebreak

\section{Begrüßung}
Als Protokollführer wird \protokoller~bestimmt.\\
Der Sprecher des Fachschaftsrates \fsiPresident~eröffnet die Vollversammlung um 13:15 Uhr.
\\\\

\textbf{Es wird festgestellt, dass die Vollversammlung beschlussfähig ist.}
\linebreak

\subsection{Anwesende}
siehe Anlagen.

\section{Satzungsänderung}
Es wird über die Satzungsänderung der Fachschaft Informatik abgestimmt.
\\\\
Es wird der Antrag gestellt, die Satzungsänderung der Fachschaftaft Informatik zu genehmigen.

\begin{center}
	\textbf{Der Antrag wird mit fünfundfünfzig Ja Stimmen, null Nein Stimmen und acht Enthaltungen angenommen.}
\end{center}
\section{Finanzen}
Der stellv. Kassenwart Philipp Dippel teilt mit, dass die Möglichkeiten zum Verkauf und zur Anschaffung weiterhin eingeschränkt sind. Der Server wurde genehmigt und wird angeschafft sobald der Jahresabschluss durch den AStA abgeschlossen wurde. Des weiteren erfolgen weitere sonstige Anschaffungen. 
Der neue Drucker kann nur vom Fachschaftsrat genutzt werden, da der Verkauf weiterhin untersagt wird.
\begin{center}
	\textbf{Der derzeitige Kontostand beträgt 4000€.}
\end{center}
Der genaue Betrag konnte durch den AStA nicht ermittelt werden.
\section{Events}
\begin{flushleft}
	Durch die Veranstaltung Glühen ist ein Verlust von 122,39€ und durch die Veranstaltung 31C3 ein Gewinn von 65,-€ entstanden. Die Teilnehmerzahl während des Glühens war leider nur gering.
	
	Die erste NIV seit langem war nur mittelmäßig gut besucht, für die nächste NIV dieses Semester soll besser geplant werden und mehr Werbung gemacht werden.
	
	Die Notld ist gut verlaufen. Fabio Gimmilaro stellt das Event noch einmal kurz vor. Während der letzten Notld wurden die NAO´s zur Verfügung gestellt. Auch während der nächsten Notld werden die NAO´s und die neu angeschafften Drohnen dabei sein.  
	
	Georg Schäfer berichtet, dass während der GDW 2 Java Spiele und 1 C++ Spiel entstanden sind. Des weiteren wird mitgeteilt, dass auch  bei der GDW dabei waren  
	
	Für den nächsten Analogen Spieleabend haben sich noch keine Studierenden gemeldet, die Spiele mitbringen werden. wenn sich bis zum 24.05. niemand mehr meldet, fällt der Abend aus.
\end{flushleft}

\section{Ratsarbeit}
\begin{flushleft}
	Für das 100-Tage Projekt hat der Fachschaftsrat in Zusammenarbeit mit Janine Kleinbauer ein Erstigrillen veranstaltet. Es ist geplant neues Mobiliar für den Fachschaftsraum anzuschaffen. Die neuen Tische sind bereits genehmigt, für die Couches liegt noch keine Genehmigung vor. 
\end{flushleft}
\section{Sonstiges}
\begin{flushleft}
	Der Präsident des StuPa Jonas Schön stellt die kommenden StuPa-Wahlen vor. Es wird angesprochen, wie wichtig es ist, das sich genügend Kandidaten aufstellen und die neue Verteilung des StuPa nach der Satzungsänderung wird erwähnt.
	\newline
	Des weiteren wird angemerkt, dass der AStA nicht länger gewählt wird, sondern durch den StuPa eingesetzt und auch hierfür bereits Kandidaten gesucht werden. 
\end{flushleft}
  

Der Sprecher Georg Schäfer schließt die Vollversammlung um 13:40 Uhr.

\pagebreak
\section{Bestätigung des Protokolls}
Der Sprecher des Fachschaftsrates Informatik sowie der Protokollführer dieses Protokolls bestätigen mit Ihrer Unterschrift unter diesem Protokoll, dass selbiges inhaltlich korrekt ist und alle darin aufgeführten Beschlüsse so wie beschrieben vom Rat getragen werden.
\\
\vspace{1.5cm}


\vspace{3.5cm}
\hrulefill \hfill \hrulefill

\fsiPresident \hfill \TeXer

{\footnotesize (Sprecher)\hfill (Protokollführer)}

\end{document}